\section{The True Sage}

\begin{quotex}
Matter --- the Platonic ``other'' --- is effectively the sign of imperfection of the act and, as it exists, is connected to that fundamental ``injustice'' about which Anaximander, Parmenides, and Empedocles speak --- and is, essentially, evil. There is no other evil beyond necessity, of which matter, brute existence, is the evidence; and as long as the ``I'' is not able to avoid the experience of matter --- of this ``other'' against and beyond him --- he will be imperfect: evil, impure, irrational.
\flright{\textsc{Julius Evola}, \emph{The Individual and the Becoming of the World}}
\end{quotex}


The fundamental difference between the Eastern way --- as in Hinduism, Buddhism, Islam --- and the Western way is the understanding of Will. The Eastern way understands the world as illusion and so the path to escape it is world rejection. The Western way understands the World as the reflection of Spirit, just as the light of the Moon is the reflection of the light from the Sun.

The World, then, is not illusion, but is rather the creation of Spirit. The Spirit is Will, Unity, Logos; on its own (that is, separated and ignorant of the Divine source), the world would is blind and deterministic. Both the divine principle and world soul have the same goal: deification or theosis of all that exists. The world soul is blind and passive and does not know what to strive for; it strives to attain unity unconsciously. The human being can participate in the freedom of the Spirit and hence be the agent for the goal of the world process.

In The \textit{Justification of the Good}, Solovyov describes the True Sage.

\begin{quotex}

In human life there is an opposition between that which is conformable to Ideas and in harmony with Reason, and that which contradicts the ideal norm. The true sage is no longer a simple hermit or a wandering monk, who has renounced life and is mildly preaching the same renunciation to others; he is one who boldly denounces the wrong and irrational things of life. Hence the end is different in the two cases. Buddha Sakyamuni peacefully dies after a meal with his disciples, while Socrates, condemned and put to prison by his fellow citizens, is sentenced by them to drink a poisoned cup. But in spite of this tragic ending, the attitude of the Greek idealist to the reality unworthy of him is not one of decisive opposition. The highest representative of humanity at this state---the philosopher---is conscious of his absolute worth in so far as he lives by pure thought in the truly existent intelligible realm of Ideas or of the all embracing rationality, and despises the false, the merely phenomenal being of the material and sensuous world.

In any case his attitude to the unworthy reality is merely one of contempt. The contempt is certainly different in kind from that characteristic of Buddhism. Buddha despises the world because everything is illusion. The very indefiniteness of this judgment, however, takes away its sting. If all is equally worthless, no one is particular is hurt by it, and if nothing but Nirvana is opposed to the bad reality, the latter may sleep in peace. For Nirvana is an absolute state and not the norm for relative states.

Now the idealist does possess such a norm and he despises and condemns the life that surrounds him not because it inevitably shares in the illusory character of everything, but because it is abnormal, irrational, opposed to the Idea. Such condemnation is no longer neutral, it has an element of defiance and demand. It is slighting to all who are bound by worldly irrationality and therefore leads to hostility, and sometimes to persecution and the cup with poison.
\end{quotex}



\flrightit{Posted on 2010-04-29 by Cologero}

\begin{center}* * *\end{center}

\begin{footnotesize}
\begin{sffamily}
\texttt{Comprehensor on 2010-04-30 at 08:42 said:}

I don't agree that there is an Eastern-Western divide, but rather it is a case of an opposition between tradition and modernism. Clearly in the Hermetic and Christian (especially biblical) texts we see the same views expressed of the world and worldly knowledge as in the East. Both Eastern and Western traditions value detachment or dispassion, which is freedom from, and control over, the passions. However, I would not call that a rejection of the world so much as discernment or putting everything in its place. For both Eastern and Western traditions encourage scientific knowledge, as it is a limited or specific mode of knowledge, but still knowledge. In modern Westernism on the other hand there is a complete negation of the Spirit. Here exists a profane culture which has nothing whatsoever to do with spiritual wisdom and does not see the world as a manifestation of the spirit.

\hfill

\texttt{Comprehensor on 2010-04-30 at 10:00 said:}

Whereas a liberated Hindu might say ``All is Brahma'' a Jacob Boehme would write, ``Thou must know that this world in its innermost unfolds its properties and powers in union with the heaven aloft above us; and so there is one Heart, one Being, one Will, one God, all in all. The outermost moving of this world cannot comprehend the outermost moving of heaven aloft above this world, for they are one to the other as life and death, or as a man and a stone are one to the other. There is a strong firmament dividing the outermost of this world from the outermost of the upper heaven; and that firmament is Death, which rules and reigns everywhere in the outermost in this world, and sets a great gulf between them.'' That one must reject the shell for the kernel is true in every orthodox tradition which is rooted in metaphysics. It therefore must be stressed that the Spirit isn't to be comprehended in the outer worlds but only through the narrow gate, as Boehme warns, ``O ye theologists, the spirit here opens a door and gate for you! If you will not now see and feed your sheep and lambs on a green meadow, instead of a dry, parched heath, you must be accountable for it before the severe, earnest and wrathful judgement of God; therefore look to it.''

\hfill

\texttt{Cologero on 2010-04-30 at 12:51 said:}

True, at the most fundamental level there is no ``divide'' in Tradition. On the other hand, different traditions have arisen to accomodate different peoples at different times. Since there is a sufficient reason for everything, it is legitimate to ask why one form rather than another is appropriate for a given race or nation.

\hfill

\texttt{Matt on 2010-04-30 at 14:04 said:}

I've been thinking about the relationship between Neo-Platonism and Gnosticism. Even though Plotinus wrote against Gnosticism, with the similarity of their view of matter as privation and non-being (therefore evil), I am beginning to wonder if once can really say there is a legitimate difference between the two. I suppose the possible difference is the dualism of much of Gnosticism, that there are two Wills, and one created created matter in opposition to the other one. No?

It also brings up the question of Evola's supposed general opposition to Gnosticism (though I have yet to read something from Evola himself that is a critique of that doctrine). Though again, I suppose it could be because of the dualism.

\hfill

\texttt{Cologero on 2010-04-30 at 18:54 said:}

Solovyov was quite familiar with Boehme. It is not at all a question about mistaking the outer world for Spirit.

What he is getting at is how to understand the world, either actively or passively. How does the appearance of the world arise?

Since it is the result of the Will, to become passive relative to the world process is to allow it to arise ``spontaneously'', in Evola's term. Solovyov, on the other hand, recognizes the task to create the world appearance with a Will that is conscious and free. As the world process is struggling to achieve union, it fails without the activity of the Will conscious of its inner unity.

This will be made clear in future posts.

\hfill

\texttt{Comprehensor on 2010-05-01 at 07:38 said:}

I don't see the difference between Buddha and Socrates. Did not Buddha boldly denounce the wrong and irrational? The only difference is that Buddhism was originally a monastic order. I don't think we can use this example as a basis for serious comparative study of different traditions or adaptations thereof.

\hfill

\texttt{Cologero on 2010-05-01 at 09:57 said:}

It is not a question of a so-called ``comparative'' study. It is about the differences in the inner states of mind associated with a given teaching, even if they are subtle. Followers of different traditions, and cultures based on those traditions, show striking differences from each other. It is hardly unreasonable to relate those outer differences to the different attitudes toward the world induced by or implicated in their respective spiritual paths.

\hfill

\texttt{Comprehensor on 2010-05-01 at 12:00 said:}

Again, these differences (of inner states and attitudes) remain to be seen, especially when based on a false premise.

\hfill

\texttt{larissa on 2010-05-17 at 18:50 said:}

``The contempt is certainly different in kind from that characteristic of Buddhism. Buddha despises the world because everything is illusion. The very indefiniteness of this judgment, however, takes away its sting.'

I suggest you make a serious study of Buddhism before making such generalized suggestions. Evola's ``Doctrine of Awakening'' is a good start. Also Rhys Davids and Coomaramaswamy are excellent as well. Buddha was merely trying to reintroduce the ``vigor'' that was dying out in the old religion--through corruptions of various sorts --which is why when one studies Buddhism carefully one realizes that the Buddha taught that which is hardly different from what Hinduism taught... just tried to reinvigorate the old ideal that were being corrupted in Hinduism... .Unfortunately a self-correcting system that was entirely native born was disrupted of course with the Islamic invasions and Hindus turning ultra--orthodox to survive a foreign despotic system imposed upon them by force.

\end{sffamily}
\end{footnotesize}

